\section{\ref{sec:motorlearning}장. 운동 학습}

최소 제곱 규칙과 개별 오차 전선의 이점은 정리이다. 오차 전선 회로의 구조, 겹칠 때의 약화, 지연에 맞춘 적격 흔적, 후효과와 기술의 자발적 복귀는 측정되었다. 회로 전체가 지도 학습으로 동작하며 내부 모델을 만든다는 것은 유력한 가설이다. 두 과정 모델의 수치는 이 책의 모델에 따른 계산이다.

{\footnotesize
\setlength{\tabcolsep}{3pt}\renewcommand{\arraystretch}{1.15}
\begin{longtable}{>{\raggedright}p{0.5cm}>{\raggedright}p{4.0cm}>{\raggedright}p{5.6cm}>{\raggedright}p{2.3cm}>{\raggedright\arraybackslash}p{1.7cm}}
\toprule
No. & 명제 & 실험과 측정한 것 & 출처 & 상태 \\
\midrule\endfirsthead
\toprule
No. & 명제 & 실험과 측정한 것 & 출처 & 상태 \\
\midrule\endhead
1 & 규칙~\eqref{eq:lms}는 평균적으로 오차 제곱의 기울기를 따라 내려간다 & 적응 필터 이론의 결과: 입력과 출력 오차에 비례하는 보정은 평균 제곱 오차의 확률적 경사 하강이다. & \cite{WidrowHoff1960} & 이론 \\
2 & 출력마다 오차 전선이 있으면 속도는 출력 수와 무관하다; 하나의 공통 수로는 잡음 내는 뉴런 수에 따라 떨어진다(명제~\ref{prop:teachers}) & 선형 네트워크의 학습 곡선: 뉴런의 무작위 섭동에 의한 학습은 정확한 기울기보다 섭동받는 뉴런 수만큼 느리다. & \cite{Werfel2005} & 이론 \\
3 & 오차 전선 회로는 지도 학습으로 동작한다(명제~\ref{prop:errorwire}) & 회로 구조에서 유도한 이론. 7--10행의 실험은 이와 부합한다; 회로 전체가 바로 이렇게 동작한다는 것은 증명되지 않았다. & \cite{Marr1969, Albus1971} & 가설 \\
4 & 확장층: 작은 \nt{GLUT} 뉴런 약 $7\cdot10^{10}$개, 뇌의 나머지 전부보다 많다 & 녹인 조직 현탁액에서 핵 계수: 사람 소뇌에 뉴런 $69\cdot10^9$개, 뇌 전체는 $86\cdot10^9$개. 입체해석학: 노년 남성의 소뇌에 작은 뉴런 $101\cdot10^9$개. & \cite{Azevedo2009, Andersen1992cereb} & 측정됨 \\
5 & 입력의 차원을 높이면 원하는 관계가 선형이 된다 & 분할 수에 관한 정리: 문턱값을 가진 입력 $n$개의 가중합은 일반 위치에 있는 벡터 약 $2n$개의 임의의 표지를 구현한다. 소뇌가 바로 이것을 쓴다는 것은 3행 가설의 일부이다. & \cite{Cover1965} & 이론 \\
6 & 출력 \nt{GABA} 뉴런 약 $3\cdot10^7$개, 각각 $10^5$개 정도의 입력 & 사람 소뇌의 입체해석학: 출력 뉴런 $30.5\cdot10^6$개. 전자 현미경, 쥐: 출력 뉴런 하나에 확장층 입력 약 $1.75\cdot10^5$개. 출력 뉴런의 억제성 신경전달물질은 총설. & \cite{Andersen1992cereb, NapperHarvey1988, Ito2001} & 측정됨 \\
7 & 출력 뉴런마다 정확히 하나의 오차 전선, 초당 약 한 번, 매번 강력한 폭발 & 기록 총설: 각 출력 뉴런은 등반하는 \nt{GLUT} 입력 하나를 받으며, 이것이 $1$~Hz 정도의 빈도로 복합 스파이크를 일으킨다. & \cite{Ito2001} & 측정됨 \\
8 & 폭발의 확률은 움직임 오차의 방향에 따라 다르다 & 원숭이, 소뇌의 안구 운동 부위, 도약 안구 운동 적응: 시각 오차의 방향이 복합 스파이크의 확률을, 크기가 그 시각을 정한다; 폭발은 다음 시도의 단순 스파이크를 바꾸고 움직임을 세포의 선호 오차 방향으로 옮긴다. & \cite{Herzfeld2018} & 측정됨 \\
9 & 오차 폭발이 활성 입력과 겹치면 그 입력이 약해진다($-\eta e_i x_j$) & 실험 총설: 확장층 입력과 오차 전선을 함께 자극하면 그 입력이 장기적으로 약해진다; 한 입력만 자극하면 그렇지 않다. & \cite{Ito2001} & 측정됨 \\
10 & 흔적의 길이는 오차 지연에 맞춰져 있다: 지연이 $100\ms$이면 $100\ms$ 앞서 활성이던 입력이 가장 강하게 약해진다 & 절편과 살아 있는 생쥐: 안구 운동 학습을 맡는 부위에서 약화는 입력과 오차 사이 간격 약 $120\ms$에 좁게 맞춰져 있다. 이 과제에서 오차 신호가 오는 데 걸리는 시간이다; 다른 부위에서는 세포마다 다른 간격에 맞춰져 있다. & \cite{Suvrathan2016} & 측정됨 \\
11 & 회로는 몸의 내부 모델을 학습하는 적응 필터~\eqref{eq:adaptivefilter}이다 & 회로 구조에서 유도한 모델. 학습된 필터를 몸의 전달 함수로 측정한 직접 실험은 없다. & \cite{Fujita1982} & 가설 \\
12 & 예측은 자기 움직임의 결과를 빼므로 스스로를 간지럽힐 수 없다 & 사람: 자기 접촉은 같은 남의 접촉보다 덜 간지럽게 느껴진다; 영상 장치에서 체감각 피질은 자기 접촉에 더 약하게 응답하고, 움직임이 피부에 닿을 때 소뇌는 덜 활성이다. 예측을 빼는 설명은 가설. & \cite{Blakemore1998} & 가설 \\
13 & 외부 힘이 있으면 빗나감이 줄고, 힘을 없애면 팔이 반대쪽으로 빗나간다 & 사람이 힘의 장 속에서 로봇 손잡이를 잡고 뻗는다: 궤적이 직선으로 돌아온다; 장을 없애면 후효과는 처음 왜곡의 거의 거울상이다; 훈련하지 않은 공간 영역으로도 전이된다. & \cite{ShadmehrMussaIvaldi1994} & 측정됨 \\
14 & 두 과정~\eqref{eq:tworate}이 학습한다; 반대 섭동 뒤에 기술이 스스로 돌아온다 & 사람, 힘의 장, 이어서 짧은 반대 장과 오차를 고정한 시도: 보정이 스스로 옛 값으로 돌아간다; 두 과정 모델은 이를 재현하고 한 과정 모델은 재현하지 못한다. & \cite{Smith2006} & 측정됨 \\
15 & 그림~\ref{fig:tworate}의 수치: $200$번 움직임에 $0.84$(느린 과정 $0.63$); 반대 $6$번; 빠른 과정 $-0.53$, 느린 과정 $0.46$; $40$번 움직임에 $0.37$까지 복귀 & \texttt{models/\allowbreak motor\_\allowbreak learning.py}에 따른 계산, $A_f=0.92$, $B_f=0.1$, $A_s=0.996$, $B_s=0.02$: $0.840$($0.634$와 $0.205$), 움직임 $6$번, $-0.531$과 $0.457$, $40$번째 움직임에서 최고 $0.371$. & \texttt{models/\allowbreak motor\_\allowbreak learning.py} & 모델 \\
16 & 몸이 여러 속도로 변하므로 두 과정이 필요하다 & 이론: 수명이 여럿인 교란에 대한 최적 추정은 시간 상수가 다른 여러 필터를 주며, 자발적 복귀와 빨라진 재학습을 설명한다. & \cite{Kording2007} & 가설 \\
\bottomrule
\end{longtable}}
